\emph{(i)} In $\Cl(1,3)$ with $Q(e_{0}) = 1$ and $Q(e_{i}) = -1$, $(e_{0}e_{j})^{2} = -Q(e_{0})Q(e_{j}) = 1$. For $v = e_{0} \pm e_{j}$, $Q(v) = 1 - 1 = 0$, and $\tfrac12(B_{j}v - vB_{j})$ equals $\mp v$, while $B_{j}$ commutes with $e_{k}$ for $k \notin \{0, j\}$. A unit boost whose timelike plane has null directions $e_{0} \pm e_{j}$ and which fixes the orthogonal spacelike plane is, at leading order, the approximate boost of the local Rindler horizons through the point generated by $e_{0} \pm e_{j}$ with that plane as edge (the leading-order form of the boost field in Theorem \claimref{synthesis.local-boost-temperature-is-local-regularity}). The check verifies these identities for $j = 1, 2, 3$. The labelling $B_{1}$ for the $t$-direction is that of Theorem \claimref{synthesis.collar-isotropy-selects-t-boost}, and $B_{1}$ is the plane of the $t$-line holonomy by Proposition \claimref{singularity.modular-flow-boundary-modes-holonomy-universal}(i).

\emph{(ii)} By (R$_{\mathrm{loc}}$) the modes of the horizon generated by $e_{0} + e_{1}$ are KMS at rapidity period $2\pi$ for its unit boost. By (A) the holonomy flow of the slice $r$ is that boost run at rate $f(r)$, so the same state is KMS for the holonomy flow parametrised by $f$ at inverse temperature $2\pi/f(r)$ (Theorem \claimref{singularity.kms-uniqueness-factor} on each window, with the reparametrisation clause of rate rigidity: a constant rescaling of the flow rescales $\beta$ inversely). With $f = \rho$ this is $2\pi \cdot 2r^{2} = 4\pi r^{2}$; at the horizon, $r = 1$, it is $4\pi = 2\pi/\kappa$ with $\kappa = \tfrac12$, the Hawking value. With $f = \rho_{\mathrm{H}}$ it is $6\pi\sqrt{1 - r}/S(r)$. The series of $S$ follows from $\tau(r) = \int_{0}^{r}\sqrt{x/(1 - x)}\,dx = \tfrac23 r^{3/2}S(r)$: expanding $(1 - x)^{-1/2} = \sum_{k}\binom{2k}{k}4^{-k}x^{k}$ and integrating term by term gives $s_{k} = \binom{2k}{k}4^{-k}\cdot\tfrac{3}{2k + 3}$, equivalently $(k + \tfrac32)s_{k} = \tfrac32\binom{2k}{k}4^{-k}$, which the check verifies to order five. Multiplying by the series of $(1 - r)^{-1/2}$ and inverting exactly gives $\rho_{\mathrm{H}} = \tfrac13 + \tfrac{4}{15}r + \cdots$ (the value stored in the cited proposition) and $\beta_{\mathrm{H}}/(6\pi) = 1 - \tfrac45 r + O(r^{2})$. The curvature scale $2M/r_{\mathrm{phys}}^{3}$ with $r_{\mathrm{phys}} = 2Mr$ equals $1/(4M^{2}r^{3})$, which is $1$ at $r_{c} = (2M)^{-2/3}$; for $2M = 1000$, $r_{c} = 1/100$ and the leading term is $6\pi \cdot \tfrac{124}{125}$. Since $\rho_{\mathrm{H}} \to \tfrac13 = |p_{1}|$ as $r \to 0$, $\beta_{\mathrm{H}} \to 6\pi$. The rapidity-period-$2\pi$ state on every slice is, under the slice-product hypothesis (C), the state of clause (iv) of Theorem \claimref{synthesis.collar-kms-extension-three-flows}, and under hypothesis (E) the state $\omega_{P}$ of Theorem \claimref{synthesis.collar-kms-extension-completion-locus}(ii); both are stated there with these inverse temperatures. The clause is used at leading order in the patch size, where the rapidity period is $2\pi$ for every curvature (Theorem \claimref{synthesis.local-boost-temperature-is-local-regularity}); near the exceptional set the patch must be small compared with the Kasner curvature radius, of order $\tau$.

\emph{(iii)} Theorem \claimref{synthesis.collar-kms-extension-three-flows}(ii) gives the Killing-extended state rapidity period $2\pi/r^{2}$ on the slice $r$. For $0 < r < 1$, $1/r^{2} > 1$, so the period is not $2\pi$ and the state violates (R$_{\mathrm{loc}}$) for the horizon generated by $e_{0} + e_{1}$ at every point of that slice; it agrees with the clause only on the horizon slice $r = 1$. Now take the one-mode-set clause of (C$_{\mathrm{loc}}$): the hole's horizon modes are the modes Premises (C1) and (C2) place at the cutoff. Premise (R) makes the horizon modes KMS for the Killing flow at $\beta = 2\pi/\kappa = 4\pi$ (units $2M = 1$); under (A) the Killing flow acts on the modes at the cutoff slice $r_{c}$ as the boost run at rapidity rate $\rho(r_{c}) = 1/(2r_{c}^{2})$, so the same state has rapidity period $4\pi\rho(r_{c}) = 2\pi/r_{c}^{2}$ there, which is the Killing extension at $r_{c}$. Since $r_{c} = (2M)^{-2/3} < 1$ for $M > \tfrac12$, this is not $2\pi$, while (R$_{\mathrm{loc}}$) at a cutoff point asks $2\pi$. So (R) at the cutoff placement and (R$_{\mathrm{loc}}$) at collar points cannot both hold on one mode set under (A). This is the obstruction found at the bifurcation surface in the shared-mode result filed alongside (synthesis.r-loc-needs-sheet-labelled-modes, consequence (a)), here between two placements of one mode set rather than two sheets of one edge. The check evaluates $1/r^{2}$ at $r = \tfrac12, \tfrac1{10}, \tfrac1{100}$ and at $r = 1$, and the cutoff period $2\pi/r_{c}^{2} = 2\pi \cdot 10^{4}$ for $2M = 1000$.

\emph{(iv)} This is the collar instance of clauses (ii) and (iv) of the shared-mode result filed alongside (synthesis.r-loc-needs-sheet-labelled-modes); the argument is short and is repeated. Read by the horizon generated by $e_{0} - e_{1}$ (same edge), whose future boost is $-B_{1}$, a pair on which the first horizon acts by $\operatorname{Ad}(e^{JsK})$ is acted on by $\operatorname{Ad}(e^{-JsK})$, $J = e_{2}e_{3}$ the completed pair's complex unit (\claimref{synthesis.transverse-bivector-is-the-complex-unit-of-the-radial-atom}); by KMS uniqueness on the factor $\Mat_{2}(\C)$ over $\mathrm{span}\{1, J\}$ the two period-$2\pi$ densities are $e^{-2\pi K}$ and $e^{2\pi K}$ up to normalisation, equal only if $K$ is a scalar, that is at zero weight (occupations $1/(1 + e^{2\pi w})$ and $1/(1 + e^{-2\pi w})$; the check shows $\tfrac16 \neq \tfrac56$ at Boltzmann symbol $\tfrac15$). On the frame atom, with $K_{j} = \tfrac12 B_{j}$ and $B_{j}^{2} = 1$, the period-$2\pi$ density for $K_{j}$ is proportional to $e^{-\pi B_{j}} = \cosh\pi - (\sinh\pi)B_{j}$; a common density would need $\cosh\pi - (\sinh\pi)B_{1} = c(\cosh\pi - (\sinh\pi)B_{2})$, impossible because $1$, $B_{1}$, $B_{2}$ are linearly independent (exact rank three in the check) and $\sinh\pi \neq 0$. This agrees with Proposition \claimref{singularity.modular-flow-boundary-modes-holonomy-universal}(iii) and Theorem \claimref{synthesis.collar-isotropy-selects-t-boost}, which reach the frame-atom statement from non-commutativity and isotropy.

\emph{Conclusion.} (ii) gives the value $6\pi$ from the clause at each point, under (S$_{1}$) and (A), so it requires no extension of the horizon's state; the values coincide with those of the rapidity-flow state of Theorem \claimref{synthesis.collar-kms-extension-three-flows}(iv), and the new content is the selection; (iii) shows that the Killing extension, which (R) gives the cutoff-placed modes on one mode set, violates the clause, a conflict between the two premises under (A); whether the period-$2\pi$ state of (ii) is also KMS for one global holonomy flow is the content of the cited theorems: yes under (E), no under (C), and (ii) does not use it.
